\section{第~\ref{sec:debugging}~章　マシンのデバッグ}

電極に何が聞こえるか、活動の低次元性、硬い配列によるインターフェースの性能、脳とデコーダの共同学習、刺激による視覚上の点は、査読付きの実験で測定されている。データ量の見積もりは本章の算術である。ヒトに埋め込んだNeuralinkの装置については、確認できた限り、データを公表しているのは主に同社自身である。

{\footnotesize
\setlength{\tabcolsep}{3pt}\renewcommand{\arraystretch}{1.15}
\begin{longtable}{>{\raggedright}p{0.5cm}>{\raggedright}p{4.0cm}>{\raggedright}p{5.6cm}>{\raggedright}p{2.3cm}>{\raggedright\arraybackslash}p{1.7cm}}
\toprule
№ & 主張 & 実験と測定内容 & 出典 & 状態 \\
\midrule\endfirsthead
\toprule
№ & 主張 & 実験と測定内容 & 出典 & 状態 \\
\midrule\endhead
1 & 電極には半径約100マイクロメートル以内のニューロン、ふつう1個から数個のスパイクが聞こえ、波形で区別できる & ラットCA1野、細胞内と細胞外の同時記録：細胞外スパイクの波形は細胞内スパイクの幅と振幅を反映する。見積もりではテトロードは$60$--$100$個のニューロンを区別できるはずだが、実際に分離できるのは約6個。 & \cite{Henze2000extra} & 測定済み \\
2 & 脳には$8.6\cdot10^{10}$個のニューロンがあり、プローブが聞くのはおよそ1億分の1 & 溶解した組織の懸濁液中の核の計数：ヒトの脳のニューロンは$86\cdot10^9$個。割合は本章の算術。 & \cite{Azevedo2009} & 測定済み \\
3 & 運動皮質の数百個のニューロンの活動は10--20個の共通変数に収まる & サル運動皮質の記録の総説：集団活動は、多くのニューロンに共通する少数の潜在変数で記述できる。 & \cite{Gallego2017} & 測定済み \\
4 & 隣接ニューロンのノイズは相関しており、100個のニューロンは1000個とほぼ同じだけ運ぶ（命題~\ref{prop:pooling}） & サル視覚皮質：隣接ニューロンのノイズ相関係数は約$0.1$で、平均化の利得は100個のニューロンで飽和する。情報制限相関の理論。 & \cite{Zohary1994, MorenoBote2014} & 測定済み \\
5 & 生データ$2\cdot10^8$~bit/sに対しスパイクでは$8\cdot10^4$~bit/s~\eqref{eq:probedata} & スパイク列の容量~\eqref{eq:spikeentropy}による本章の算術。デジタル化のパラメータは実際のもの：Neuralinkのチップは$19.3$~kHz、チャネルあたり$10$~bit。 & 本章で導出；\cite{Rieke1997, Musk2019} & 理論 \\
6 & Neuralinkの最初のシステム：チップは増幅とデジタル化を行い、生データはケーブルで送られ、スパイクは外部のプログラムが検出する。チップ上の検出、密閉筐体、無線、ワイヤレス給電は、同社によればヒト用の装置のもの & 同社の論文：生データ全体をUSB-Cケーブルで送り、チップ外のプログラムがリアルタイムでスパイクを検出する。搭載した圧縮、ワイヤレス給電と伝送は臨床用装置の計画とされている。ヒトに埋め込んだ装置については同社の発表のみ。チップ上でのスパイク検出が必要だというのは、\eqref{eq:probedata}からの本章の結論。 & \cite{Musk2019}；Neuralinkの報告、査読付き論文なし & 測定済み（同社の論文）。ヒトについては同社の報告 \\
7 & $96$本のスレッドに$32$個ずつ最大$3072$個の電極。スレッドはロボットが血管を避けて挿入する & 同社の論文：$96$本のスレッドに$32$個ずつ$3072$個の電極。ロボットは毎分$6$本（$192$電極）を挿入する。配列を長期埋め込みしたラットで最大$70\,\%$の電極がスパイクを聞き取る。 & \cite{Musk2019} & 測定済み（同社の論文） \\
8 & ヒトでは$64$本のスレッドで約1000チャネル。一部のスレッドがずれた。カーソルは約$10$~bit/s & 同社の発表のみ。PubMedの検索ではヒトでの結果を示す査読付き論文は見つからなかった。本文中の但し書きと一致する。 & Neuralinkの報告。査読付き論文なし & 測定済み（同社の報告） \\
9 & 100電極の配列によるインターフェースがカーソルを操作する & 麻痺のある人、一次運動皮質に$96$電極：受傷3年後でも腕を動かす意図がスパイクを変える。デコーダは「ニューラルカーソル」（メール、テレビ）、義手と ロボットアームの操作を可能にする。 & \cite{Hochberg2006} & 測定済み \\
10 & 「想像上の手」による書字：毎分約$90$文字、$8$~bit/s未満 & 手が麻痺した人が手書きを試み、再帰型デコーダを使用：リアルタイムで毎分$90$文字、精度$94.1\,\%$、自動修正後は$99\,\%$超。ビットでの見積もりは本章の算術。 & \cite{Willett2021} & 測定済み \\
11 & 話そうとする試みが毎分約$60$語の速さでテキストに変換される & 明瞭な発話を失った人、皮質に配列：毎分$62$語。語彙$50$語で単語誤り率$9.1\,\%$、語彙$125\,000$語で$23.8\,\%$。 & \cite{Willett2023} & 測定済み \\
12 & インターフェースが取り出すのは容量のごく一部：1ニューロンで毎秒数十ビット、100個で数千ビット（命題~\ref{prop:probe}） & 上限~\eqref{eq:spikeentropy}は理論。8行目、10行目との比較は本章の結論。ボトルネックが配線ではなく低次元性と符号の無知であることは、直接の実験で確かめられていない。 & \cite{Rieke1997} & 理論 \\
13 & 脳とデコーダは互いから学ぶ & サルがカーソルを操作。デコーダは閉ループで適応し、ニューロンも同時に学習し直す。精度は上がり、技能は保持され、自分の腕の動きの影響を受けにくくなる。学習速度を分けるという助言は工学上の経験則であり、本章に個別の実験はない。 & \cite{Orsborn2014coad} & 測定済み \\
14 & 電極を流れる電流は、周囲のニューロンと配線を種類の区別なく興奮させる & マウスとネコの皮質、二光子カルシウムイメージング：弱い電流でも、直径数十マイクロメートルの体積内の軸索を直接興奮させることで、最大数ミリメートルに散らばるまばらなニューロンを興奮させる。つまり興奮は非選択的だが、一様ではない。 & \cite{Histed2009stim} & 測定済み \\
15 & 視覚皮質の刺激は点を灯す。同時に最大$16$点で一部の文字と物体の境界が認識できる & 盲目の人、視覚皮質に$96$電極を$6$か月：単一電極の閾値は$66.8\pm36.5$~\textmu A。電極群の同時刺激（最大$16$電極は刺激装置の上限）で閾値が下がり、区別できるパターンが生じ、一部の文字と物体の境界が認識される。 & \cite{Fernandez2021} & 測定済み \\
16 & Neuralinkの第二の方向は、視覚皮質に信号を入れる装置 & 開発についての同社の発表のみ。その動作に関する査読付きデータは見つからなかった。 & Neuralinkの報告 & 仮説 \\
17 & ブロードキャスト伝達物質の濃度は組織内の化学センサで測れる & ラット脳スライス、炭素繊維による高速サイクリックボルタンメトリー：セロトニンの放出と再取り込みを測定。物質はシナプスの外へ出る。 & \cite{Bunin1998} & 測定済み \\
\bottomrule
\end{longtable}}
